\documentclass{article}
\usepackage{amsmath, amssymb, amsthm}
\usepackage{hyperref}
\usepackage{geometry}
\geometry{margin=1in}

\title{Erd\H{o}s Problem \#930}
\date{Last updated: 2025-08-31}
\author{Source: erdosproblems.com}

\begin{document}
\maketitle

\noindent\textbf{Status:} OPEN \quad \textbf{No prize}

\noindent\textbf{Tags:} number theory

\medskip

\noindent\textbf{Problem Statement:}

Is it true that, for every $r$, there is a $k$ such that if $I_1,\ldots,I_r$ are disjoint intervals of consecutive integers, all of length at least $k$, then\[\prod_{1\leq i\leq r}\prod_{m\in I_i}m\]is not a perfect power?




Erd\H{o}s and Selfridge \cite{ErSe75} proved that the product of consecutive integers is never a power (establishing the case $r=1$). The condition that the intervals be large in terms of $r$ is necessary for $r=2$ - see the constructions in [363] . See also [363] for the case of squares.




References

[ErSe75] Erd\H{o}s, P. and Selfridge, J. L., The product of consecutive integers is never a power . Illinois J. Math. (1975), 292-301.


\bigskip
\noindent\small{Source: \url{https://www.erdosproblems.com/930}}
\end{document}
